% English translation of chapter7.tex lines 611--711.
% Source: Methods of Algebra, Volume 2, commit
% 9a5803ff77dd3257484cb177f851a73770a59dd3 (CC BY 4.0).
% stable-unit-id: o014.aljabr2.chapter7.closed-monoidal-categories
% Coverage: all of Section 7.3 on closed monoidal categories; ends immediately
% before Section 7.4 at line 712.

% segment-id: o014.aljabr2.chapter7.closed-monoidal-categories.h001
\section{Closed Monoidal Categories}\label{sec:closed-monoidal}
\hypertarget{o014.aljabr2.chapter7.closed-monoidal-categories}{ }

% segment-id: o014.aljabr2.chapter7.closed-monoidal-categories.p001
Let $\Bbbk$ be a commutative ring. Then $\Bbbk\dcate{Mod}$ is a symmetric
monoidal category under $\otimes:=\otimes_{\Bbbk}$. It has the following
closedness properties.
% segment-id: o014.aljabr2.chapter7.closed-monoidal-categories.l001
\begin{itemize}
% segment-id: o014.aljabr2.chapter7.closed-monoidal-categories.i001
	\item It is enriched over itself: the set of homomorphisms $\Hom(X,Y)$
	naturally has the structure of a $\Bbbk$-module;
% segment-id: o014.aljabr2.chapter7.closed-monoidal-categories.i002
	\item the basic properties of the tensor product give a canonical
	isomorphism in $\Bbbk\dcate{Mod}$
% segment-id: o014.aljabr2.chapter7.closed-monoidal-categories.q001
	\[ \Hom(X \otimes Y, Z) \simeq \Hom(X, \Hom(Y, Z)). \]
\end{itemize}

% segment-id: o014.aljabr2.chapter7.closed-monoidal-categories.p002
Extending these properties to general monoidal categories leads to the
following concept.

% segment-id: o014.aljabr2.chapter7.closed-monoidal-categories.d001
\begin{definition}\label{def:closed-monoidal-cat}
% segment-id: o014.aljabr2.chapter7.closed-monoidal-categories.x001
	\index{monoidal category!closed}
	Let $\mathcal{C}$ be a monoidal category. The category $\mathcal{C}$ is
	called \emph{right closed} (or \emph{left closed}) if, for every object
	$Y$, the functor $(\cdot)\otimes Y:\mathcal{C}\to\mathcal{C}$ (or
	$Y\otimes(\cdot):\mathcal{C}\to\mathcal{C}$) is equipped with a specified
	right adjoint. A monoidal category that is both left closed and right
	closed is called a \emph{closed monoidal category}.
\end{definition}

% segment-id: o014.aljabr2.chapter7.closed-monoidal-categories.p003
All monoidal categories considered below are braided monoidal categories, so
we shall no longer distinguish left closed from right closed. In this
situation, the right adjoint of $(\cdot)\otimes Y$ is denoted by
$[Y,\cdot]:\mathcal{C}\to\mathcal{C}$.

% segment-id: o014.aljabr2.chapter7.closed-monoidal-categories.p004
Recall that if a category $\mathcal{C}$ has finite products, then
$\mathcal{C}$ is a symmetric monoidal category under $\times$, with the
terminal object as its unit $\munit$ (Example~\ref{eg:Cartesius-monoidal}).

% segment-id: o014.aljabr2.chapter7.closed-monoidal-categories.d002
\begin{definition}\label{def:Cartesian-closed}
% segment-id: o014.aljabr2.chapter7.closed-monoidal-categories.x002
	\index{category!Cartesian closed}
	Let $\mathcal{C}$ be a category with finite products. If
	$(\mathcal{C},\times)$ is a closed monoidal category, then $\mathcal{C}$ is
	called a \emph{Cartesian closed category}.
\end{definition}

% segment-id: o014.aljabr2.chapter7.closed-monoidal-categories.p005
More generally, for any two pairs of adjoint functors $(F,G)$ and $(F',G')$
between categories $\mathcal{C}_1$ and $\mathcal{C}_2$, every
$\varphi:F\to F'$ naturally induces $\psi:G'\to G$; conversely, every
$\psi:G'\to G$ naturally induces $\varphi:F\to F'$. In either case, the
induced morphism is characterized by the commutative diagram
% segment-id: o014.aljabr2.chapter7.closed-monoidal-categories.g001
\[\begin{tikzcd}
	\Hom(F' X, Y) \arrow[r, "\sim"] \arrow[d, "{(\varphi_X)^*}"'] & \Hom(X, G'Y) \arrow[d, "{(\psi_Y)_*}"] \\
	\Hom(FX, Y) \arrow[r, "\sim"'] & \Hom(X, GY)
\end{tikzcd}\]
This is simply an application of the Yoneda lemma. The reader may also try to
write down these induced morphisms using the unit and counit of the adjoint
pairs.

% segment-id: o014.aljabr2.chapter7.closed-monoidal-categories.p006
As an application, every morphism $Y\to Y'$ in a closed monoidal category
induces $[Y',\cdot]\to[Y,\cdot]$. Thus, being a closed monoidal category is
equivalent to the existence of a bifunctor
$[\cdot,\cdot]:\mathcal{C}^{\opp}\times\mathcal{C}\to\mathcal{C}$ and a
family of canonical bijections
% segment-id: o014.aljabr2.chapter7.closed-monoidal-categories.e001
\begin{equation}\label{eqn:closed-monoidal-adj0}
	\Hom(X \otimes Y, Z) \simeq \Hom(X, [Y, Z]),
\end{equation}
functorial in all three variables.

% segment-id: o014.aljabr2.chapter7.closed-monoidal-categories.p007
The bifunctor $[\cdot,\cdot]$ is also called the \emph{internal $\Hom$} of the
closed monoidal category $\mathcal{C}$. The definition also yields the
following conclusions.
% segment-id: o014.aljabr2.chapter7.closed-monoidal-categories.l002
\begin{itemize}
% segment-id: o014.aljabr2.chapter7.closed-monoidal-categories.i003
	\item From
	$\Hom(X,Z)\simeq\Hom(X\otimes\munit,Z)\simeq\Hom(X,[\munit,Z])$ and the
	Yoneda lemma, one obtains $Z\simeq[\munit,Z]$;
% segment-id: o014.aljabr2.chapter7.closed-monoidal-categories.i004
	\item the unit morphism of the adjoint pair gives
	$\mathrm{coev}_{X,Y}:X\to[Y,X\otimes Y]$, while the counit morphism gives
	$\mathrm{ev}_{Y,X}:[Y,X]\otimes Y\to X$;
% segment-id: o014.aljabr2.chapter7.closed-monoidal-categories.i005
	\item from the composite
% segment-id: o014.aljabr2.chapter7.closed-monoidal-categories.q002
	\[ [Y, Z] \otimes ([X, Y] \otimes X) \xrightarrow{\identity \otimes \mathrm{ev}_{X, Y}} [Y, Z] \otimes Y \xrightarrow{\mathrm{ev}_{Y, Z}} Z \]
	and the adjunction together with the associativity constraint, one obtains
% segment-id: o014.aljabr2.chapter7.closed-monoidal-categories.q003
	\[ [Y, Z] \otimes [X, Y] \to [X, Z]; \]
% segment-id: o014.aljabr2.chapter7.closed-monoidal-categories.i006
	\item taking $\mathrm{coev}_{\munit,X}$ gives a morphism
	$\munit\to[X,X]$.
\end{itemize}

% segment-id: o014.aljabr2.chapter7.closed-monoidal-categories.p008
The term internal $\Hom$ and the preceding properties make clear that a closed
monoidal category is enriched over itself. We first show how to recover the
ordinary $\Hom$ from the internal $\Hom$.

% segment-id: o014.aljabr2.chapter7.closed-monoidal-categories.pr001
\begin{proposition}
	Let $\mathcal{C}$ be a closed monoidal category. There is a family of
	canonical bijections
% segment-id: o014.aljabr2.chapter7.closed-monoidal-categories.q004
	\[ \Hom(X, Y) \simeq \Hom(\munit, [X, Y]), \quad X, Y \in \Obj(\mathcal{C}). \]
\end{proposition}

% segment-id: o014.aljabr2.chapter7.closed-monoidal-categories.pf001
\begin{proof}
	By \eqref{eqn:closed-monoidal-adj0},
	$\Hom(\munit,[X,Y])\simeq\Hom(\munit\otimes X,Y)\simeq\Hom(X,Y)$.
\end{proof}

% segment-id: o014.aljabr2.chapter7.closed-monoidal-categories.p009
To show that $\mathcal{C}$ is enriched over itself, one must still prove that
$[Y,Z]\otimes[X,Z]\to[Y,Z]$ and $\munit\to[X,X]$ satisfy axioms such as
associativity. This can be checked using the standard properties of the unit
and counit. Since the details are rather tedious, the proof is left as an
exercise for this chapter.

% segment-id: o014.aljabr2.chapter7.closed-monoidal-categories.p010
We next show that the adjunction \eqref{eqn:closed-monoidal-adj0} is itself
automatically internalized in $\mathcal{C}$.

% segment-id: o014.aljabr2.chapter7.closed-monoidal-categories.pr002
\begin{proposition}\label{prop:closed-monoidal-adj}
	Let $\mathcal{C}$ be a closed monoidal category. There is a family of
	canonical isomorphisms
% segment-id: o014.aljabr2.chapter7.closed-monoidal-categories.q005
	\[ [X \otimes Y, Z] \simeq [X, [Y, Z]]; \]
	more precisely, this is an isomorphism between functors from
	$\mathcal{C}^{\opp}\times\mathcal{C}^{\opp}\times\mathcal{C}$ to
	$\mathcal{C}$.
\end{proposition}

% segment-id: o014.aljabr2.chapter7.closed-monoidal-categories.pf002
\begin{proof}
	Fix $X$, $Y$, and $Z$. For every object $T$,
	\eqref{eqn:closed-monoidal-adj0} and the associativity constraint of
	$\mathcal{C}$ give natural bijections
% segment-id: o014.aljabr2.chapter7.closed-monoidal-categories.g002
	\begin{multline*}
		\Hom(T, [X \otimes Y, Z]) \simeq \Hom(T \otimes (X \otimes Y), Z) \simeq \Hom((T \otimes X) \otimes Y, Z) \\
		\simeq \Hom(T \otimes X, [Y, Z]) \simeq \Hom(T, [X, [Y, Z]]).
	\end{multline*}
	Since $T$ is arbitrary, the Yoneda lemma gives the desired isomorphism.
\end{proof}

% segment-id: o014.aljabr2.chapter7.closed-monoidal-categories.p011
The following basic examples involve the symmetric monoidal categories
introduced in \S\ref{sec:alg-in-monoidal-cat}.

% segment-id: o014.aljabr2.chapter7.closed-monoidal-categories.l003
\begin{itemize}
% segment-id: o014.aljabr2.chapter7.closed-monoidal-categories.i007
	\item The category of sets $\cate{Set}$ is Cartesian closed: take $[X,Y]$
	to be the set of maps
	$Y^X=\left\{\text{maps}\;X\to Y\right\}$. Then
	\eqref{eqn:closed-monoidal-adj0}, or its internalized version in
	Proposition~\ref{prop:closed-monoidal-adj}, gives the natural bijection
	$Z^{X\times Y}\xrightarrow{1:1}(Z^X)^Y$. In computer science, such a
	bijection, or its generalization to an arbitrary closed monoidal category,
	is often called \emph{currying}.

% segment-id: o014.aljabr2.chapter7.closed-monoidal-categories.i008
	\item All small categories and the functors between them form the category
	$\cate{Cat}$. Its products are the products of categories
	$\mathcal{C}_1\times\cdots\times\mathcal{C}_2$, and its empty product is
	the category $\mathbf{1}$. The category $\cate{Cat}$ is Cartesian closed;
	this is equivalent to the following simple statement: specifying a
	bifunctor $\mathcal{A}\times\mathcal{B}\to\mathcal{C}$ is equivalent to
	specifying a functor $\mathcal{A}\to\mathcal{B}^{\mathcal{C}}$, and is
	also equivalent to specifying $\mathcal{B}\to\mathcal{C}^{\mathcal{A}}$.

% segment-id: o014.aljabr2.chapter7.closed-monoidal-categories.i009
	\item Although the category of topological spaces $\cate{Top}$ has many
	good properties and is a symmetric monoidal category under the product
	$\times$, it is not Cartesian closed; see the discussion at the end of
	\cite[\S 1.5]{Kel05}. Since mapping spaces $X^Y$ occur everywhere in
	topology, the cost of giving up closedness is too high. A convenient
	replacement is the category of compactly generated Hausdorff
	spaces\footnote{A compactly generated Hausdorff space is a Hausdorff space
	$X$ with the following property: if the intersection of a subset
	$A\subset X$ with every compact subset $K\subset X$ is closed, then $A$
	is closed.} $\cate{CGHaus}$. It is a full subcategory of $\cate{Top}$ and
	is itself Cartesian closed. One can prove that there is an adjoint pair
% segment-id: o014.aljabr2.chapter7.closed-monoidal-categories.g003
	$\begin{tikzcd}
		\cate{CGHaus} \arrow[shift left, r, "\text{inclusion functor}"] & \cate{Top} \arrow[shift left, l, "k"]
	\end{tikzcd}$,
	where the inclusion functor preserves $\varinjlim$ but not products.

% segment-id: o014.aljabr2.chapter7.closed-monoidal-categories.i010
	\item Let $G$ be a group and $\Bbbk$ a commutative ring. The symmetric
	monoidal category $G\dcate{Mod}$ introduced in
	Proposition~\ref{prop:Gmod-monoidal} is closed: its internal $\Hom$ is
	precisely the $\Hom$ module introduced in \S\ref{sec:G-mod}, and the
	required adjunction is exactly \eqref{eqn:G-mod-tensor-Hom}.

% segment-id: o014.aljabr2.chapter7.closed-monoidal-categories.i011
	\item Again let $\Bbbk$ be a commutative ring, and set
	$\mathcal{A}:=\Bbbk\dcate{Mod}$. As seen in \S\ref{sec:dga-monoidal}, the
	category of complexes $\cate{C}(\mathcal{A})$ is a symmetric monoidal
	category. It is closed, with internal $\Hom$ given by the $\Hom$ complex.
	The adjunction required in \eqref{eqn:closed-monoidal-adj0} is the
	adjunction between $\Hom$ and $\otimes$ at the level of complexes.
	Concretely, it is a corollary of Proposition~\ref{prop:tensor-Hom-cplx} in
	the case $A=R=B=\Bbbk$. The corresponding unit and counit morphisms are
% segment-id: o014.aljabr2.chapter7.closed-monoidal-categories.g004
	\[\begin{tikzcd}[row sep=tiny]
		\mathrm{ev}_{X, Y}: \Hom^n(X, Y) \dotimes{\Bbbk} X^m \arrow[r] & Y^{n+m} \\
		f \otimes x \arrow[mapsto, r] & f(x), \\
		\mathrm{coev}_{X, Y}: X^n \arrow[r] & \Hom^n(Y, X \otimes Y) \\
		x \arrow[mapsto, r] & {[y \mapsto x \otimes y]}.
	\end{tikzcd}\]
\end{itemize}

% segment-id: o014.aljabr2.chapter7.closed-monoidal-categories.p012
The last example shows that Definition~\ref{def:Hom-cplx} of the $\Hom$
complex is entirely natural, provided that $\cate{C}(\mathcal{A})$ is given
the symmetric monoidal structure described above.
